\section{Results}
\label{sec:results}

Every number below is taken from the saved tables for \texttt{extracao=2026-03}.
The rank is test sMAPE.
Lower is better.
The skill score is higher when the model beats the seasonal naive of that grain.

\subsection{The lowest sMAPE is not the same model on every grain}
\label{sec:top3}

Table~\ref{tab:top3} lists the three lowest test-sMAPE models on each grain.
E2 has the lowest test sMAPE on the monthly panel.
TimesFM 3.0 has the lowest test sMAPE on grain A, on grain B, and on the weekly series.
A single winner across grains is not supported by this table.

\begin{table}[t]
  \caption{Three lowest test sMAPE values on each grain.
  E2 is the disagreement proxy.
  E2 is not TEM.}
  \label{tab:top3}
  \centering
  \small
  % Auto-built from freeze ranked CSVs — do not hand-edit numbers
\begin{tabular}{llrr}
\toprule
Grain & Model & sMAPE $\downarrow$ & SS\_sMAPE $\uparrow$ \\
\midrule
Monthly & E2 (proxy ablat.) & 10.94 & 0.192 \\
 & M3\_HGB\_tuned & 13.34 & 0.015 \\
 & naive seas lag12 & 13.54 & 0.000 \\
\addlinespace
Daily A (calendar) & M9\_TimesFM3 & 59.76 & 0.200 \\
 & naive seas lag7 & 74.68 & 0.000 \\
 & M6\_ARF\_regressor & 85.15 & -0.140 \\
\addlinespace
Daily B (business) & M9\_TimesFM3 & 52.39 & 0.284 \\
 & M3\_HGB\_tuned & 53.70 & 0.266 \\
 & M1\_OLS & 54.05 & 0.261 \\
\addlinespace
Weekly & M9\_TimesFM3 & 17.25 & 0.415 \\
 & M1\_Ridge\_a0.1 & 25.11 & 0.148 \\
 & M1\_OLS & 26.51 & 0.100 \\
\addlinespace
\bottomrule
\end{tabular}

\end{table}

\subsection{Monthly panel}
\label{sec:monthly}

E2 has the lowest test sMAPE on the monthly panel.
The value is about 10.9, and the skill score is about 0.19.
That is lower error than boosting alone and lower error than the lag-12 seasonal naive.

Boosting is second.
Its test sMAPE is about 13.3, and its skill score is about 0.015.
The lag-12 seasonal naive is third, at about 13.5.
Its skill score is zero by~\eqref{eq:ss}.

Selective TEM does not have the lowest monthly test sMAPE.
E1 is about 15.5, and its skill score is negative.
On this 12-month test, the energy rule has higher sMAPE than the seasonal naive.

Table~\ref{tab:monthly} continues the list to eight models.
TimesFM and the score sampler are both worse than the seasonal naive on this grain.

\begin{table}[t]
  \caption{Monthly grain: eight lowest test sMAPE values.}
  \label{tab:monthly}
  \centering
  \small
  % Monthly top-8 from metrics_monthly_test_ranked.csv
\begin{tabular}{lrrl}
\toprule
Model & sMAPE $\downarrow$ & SS\_sMAPE $\uparrow$ & Notes \\
\midrule
E2 (proxy ablat.) & 10.94 & 0.192 & proxy ablation $\neq$ TEM \\
M3\_HGB\_tuned & 13.34 & 0.015 & HistGradientBoosting \\
naive seas lag12 & 13.54 & 0.000 & seasonal ref ($SS{=}0$) \\
M8\_KRR\_rbf & 14.67 & -0.083 & Kernel Ridge \\
E1 TEM selective & 15.45 & -0.141 & real TEM selective \\
naive lag1 & 16.09 & -0.188 & lag-1 naive \\
M9\_TimesFM3 & 18.62 & -0.375 & TimesFM 3.0 zero-shot \\
E4\_ScoreGrad & 20.91 & -0.544 & score + Langevin \\
\bottomrule
\end{tabular}

\end{table}

\subsection{Grain A and grain B}
\label{sec:ablation}

TimesFM has the lowest test sMAPE on calendar-daily A.
The value is about 59.8, and the skill score is about 0.20.
The lag-7 seasonal naive is second.
Most other models have higher sMAPE than that naive, so their skill scores are negative.
Weekend and empty-slip zeros are in this score.

TimesFM has the lowest test sMAPE on business-daily B.
The value is about 52.4, and the skill score is about 0.28.
Boosting and ordinary least squares are next (Table~\ref{tab:ablation}).
Dropping weekends lowers the sMAPE level and changes the order.
The seasonal naive is inside the top three on A and outside the top three on B.

The positive-day score on A is not the official rank.
On days with $y \gt 0$, TimesFM has sMAPE about 29.9 and skill about 0.58.
Boosting has sMAPE about 34.6 and skill about 0.51.
The gap versus the all-day scores shows how much Ranking A is raised by zeros.
It is not a reason to remove those zeros from the target.

\begin{table}[t]
  \caption{Three lowest test sMAPE values on grain A and on grain B.}
  \label{tab:ablation}
  \centering
  \small
  % From metrics_ablation_A_vs_B_top3.csv
\begin{tabular}{llrr}
\toprule
Target & Model & sMAPE $\downarrow$ & SS\_sMAPE $\uparrow$ \\
\midrule
A calendar & M9\_TimesFM3 & 59.76 & 0.200 \\
A calendar & naive seas lag7 & 74.68 & 0.000 \\
A calendar & M6\_ARF\_regressor & 85.15 & -0.140 \\
B business days & M9\_TimesFM3 & 52.39 & 0.284 \\
B business days & M3\_HGB\_tuned & 53.70 & 0.266 \\
B business days & M1\_OLS & 54.05 & 0.261 \\
\bottomrule
\end{tabular}

\end{table}

\subsection{Weekly series}
\label{sec:weekly}

TimesFM has the lowest test sMAPE on the weekly series.
The value is about 17.2, and the skill score against the lag-52 naive is about 0.41.
Ridge and ordinary least squares are next.
Most other models have negative skill against that naive.
The test block has 13 weeks.
The gap describes this window.
It is not a standard error for a long-run difference.
